% P15: the separate Feller II.2 weak-law and median/symmetrization presentation.
\chapter{Feller's Symmetrization Proof of the Weak Law}
\label{sec:p15-feller-weak-law-presentation}

Feller's generalized weak law uses a median to compare the tails of a
distribution with those of its symmetrization. Its principal estimate
comes from the largest summand and the sign symmetry of the remaining
sum. This is a different source route from the first-crossing proof in
Chapter~\ref{sec:p15-weak-laws-and-triangular-arrays}.

The presentation is Feller, Volume II, second edition: the median
definition in V.2, the definitions V.(5.1)--(5.3) and inequalities
V.(5.6)--(5.10),
the integration-by-parts argument V.(6.1), and VII.7 through Theorem 1
\cite[pp.~137, 148--150, 234--236]{Feller1971ProbabilityII}.
Every premise used by the median chain is proved below. Endpoint
conventions are explicit for atomic laws.

\section{Medians, reflection and symmetrization}
\label{sec:p15-feller-median-symmetrization}

\begin{hbdefinition}{def:p15-feller-law-conventions}{Feller's distribution notation and medians}
Let \(X\) be a real random variable on
\((\Omega,\mathcal F,\mathbb P)\). Write
\(F(x)=\mathbb P(X\leq x)\) for its right-continuous CDF,
and \(F\{dx\}\) for integration against its associated Borel
probability measure. In particular
\(F(x-)=\mathbb P(X<x)\).
A median is a finite real \(m\) such that
\[
 F(m)\geq\tfrac12,\qquad F(m-)\leq\tfrac12,
\]
equivalently \(\mathbb P(X\leq m)\geq1/2\) and
\(\mathbb P(X\geq m)\geq1/2\).

Following the source, \(-F\) denotes the law of \(-X\); it is
not the pointwise negative of a CDF. For independent copies \(X,X'\)
of the same law, put
\({}^{\circ}X=X-X'\), and let \({}^{\circ}F\) be its CDF.
The operation is called symmetrization. A law is symmetric about zero
when the laws of \(X\) and \(-X\) agree.
\index{median}\index{symmetrization}
\end{hbdefinition}

\begin{hbtheorem}{thm:p15-feller-symmetrization-structure}{proved=annex:p15.feller-symmetrization-structure}{Medians and the exact reflected and symmetrized laws}
Every real probability law has a median, which need not be unique.
Centering \(X\) by a median makes zero a median.
For every real \(x\),
\[
 (-F)(x)=1-F((-x)-),\qquad
 {}^{\circ}F(x)=\int_{\mathbb R}F(x+y)\,F\{dy\}.
\]
The symmetrized law is \(F\star(-F)\), the convolution of the
two indicated laws, and is symmetric:
\[
 {}^{\circ}F(x)=1-{}^{\circ}F((-x)-).
\]
The source formulas without the left-limit sign apply at the
corresponding continuity points.
Independent copies of an entire sequence can be constructed on the
finite product of its original probability space with itself. If the
original sequence is independent, its copied differences are independent;
if it is also identically distributed, so are the differences.
Symmetrization is unchanged by any deterministic translation.
\end{hbtheorem}

\begin{hbtheorem}{thm:p15-feller-symmetrization-inequalities}{proved=annex:p15.feller-symmetrization-inequalities}{The complete tail comparisons V.(5.6)--(5.8)}
For independent identically distributed real \(X,X'\) and \(t>0\),
\[
 \mathbb P(|X-X'|>t)\leq2\mathbb P(|X|>t/2).
\]
If \(a\geq0\), \(p\in[0,1]\), and
\(\mathbb P(X\leq a)\geq p\),
\(\mathbb P(X\geq-a)\geq p\), then
\[
 \mathbb P(|X-X'|>t)\geq p\,\mathbb P(|X|>t+a).
\]
In particular, if zero is a median, the right side can be
\(\tfrac12\mathbb P(|X|>t)\).
For an arbitrary median \(m\),
\[
 \mathbb P(|X-X'|>t)
           \geq\tfrac12\mathbb P(|X|>t+|m|).
\]
For every finite real \(b\) the upper bound may also be written
\(\mathbb P(|X-X'|>2t)\leq2\mathbb P(|X-b|>t)\).
None of these statements assumes that an expectation exists.
\end{hbtheorem}

\section{Feller's largest-summand argument}
\label{sec:p15-feller-largest-summand}

\begin{hbdefinition}{def:p15-feller-largest-term}{The first largest absolute summand}
For independent symmetric real \(Y_1,\ldots,Y_n\), \(n\geq1\),
write \(S=\sum_{i=1}^nY_i\), and define
\[
 J=\min\{j:|Y_j|=\max_{1\leq i\leq n}|Y_i|\},\qquad
 M=Y_J,\qquad R=S-M.
\]
The first index resolves every tie, including the case in which all
summands vanish. Both \(M\) and \(R\) are measurable. Their roles
are a selected signed summand and the sum of the others; they are
not assumed independent.
\end{hbdefinition}

\begin{hbtheorem}{thm:p15-feller-largest-term-inequality}{proved=annex:p15.feller-largest-term-inequality}{The source inequalities V.(5.9)--(5.10), including atoms}
Under the independent symmetric-summand hypotheses of
Definition~\ref{def:p15-feller-largest-term},
the four pairs \((M,R),(-M,R),(M,-R),(-M,-R)\) have the same
law. For every \(t>0\),
\[
 \mathbb P(|S|>t)\geq\tfrac12
                    \mathbb P(\max_i|Y_i|>t),\qquad
 \mathbb P(|S|\geq t)\geq\tfrac12
                    \mathbb P(\max_i|Y_i|\geq t).
\]
If the \(Y_i\) have a common CDF \(G\), put
\(q_>(t)=\mathbb P(|Y_1|>t)\) and
\(q_\geq(t)=\mathbb P(|Y_1|\geq t)\). Then
\[
 \begin{aligned}
 \mathbb P(|S|>t)&\geq\tfrac12[1-(1-q_>(t))^n]
                         \geq\tfrac12[1-e^{-nq_>(t)}],\\
 \mathbb P(|S|\geq t)&\geq\tfrac12[1-(1-q_\geq(t))^n]
                         \geq\tfrac12[1-e^{-nq_\geq(t)}].
 \end{aligned}
\]
In particular, the exact source CDF form is valid at every \(t>0\):
\[
 \mathbb P(|S|\geq t)
       \geq\tfrac12\left(1-e^{-n[1-G(t)+G(-t)]}\right).
\]
The strict-tail formula in the preceding display is the version used
in the weak-law necessity chain.
\end{hbtheorem}

\section{The truncated moment identity}
\label{sec:p15-feller-truncated-moment}

\begin{hbtheorem}{thm:p15-feller-integration-by-parts}{proved=annex:p15.feller-integration-by-parts}{The source finite-interval Stieltjes formula}
Let \(a<b\) be finite, let \(u\) be continuously differentiable
on \([a,b]\), and let \(F\) be a CDF. Then
\[
 \int_{(a,b]}u(x)\,F\{dx\}
 =u(b)F(b)-u(a)F(a)-\int_a^b u'(x)F(x)\,dx.
\]
The first integral excludes \(a\) and includes \(b\); the last
integral is with respect to Borel Lebesgue measure. All terms are
finite. The annex supplies the source's partition argument, including
the finite summation identity and both approximation errors.
\end{hbtheorem}

\begin{hbdefinition}{def:p15-feller-tail-moment-functions}{The source tail quantity and explicit endpoint conventions}
For the common law \(F\), and a real \(t>0\), put
\[
 \begin{aligned}
 \tau(t)&=t[1-F(t)+F(-t)]
       =t[\mathbb P(X>t)+\mathbb P(X\leq-t)],\\
 \sigma(t)&=\frac1t\int_{[-t,t]}x^2\,F\{dx\},&
 \sigma_h(t)&=\frac1t\int_{(-t,t]}x^2\,F\{dx\}.
 \end{aligned}
\]
The closed version \(\sigma\) corresponds to the source's deletion
truncation \(X\mathbf1_{\{|X|\leq t\}}\); \(\sigma_h\) records
the half-open Stieltjes interval. Define the corresponding first
truncated moments by
\[
 m(t)=\int_{[-t,t]}x\,F\{dx\},\qquad
 m_h(t)=\int_{(-t,t]}x\,F\{dx\}.
\]
All four truncated integrals are finite without a moment assumption
on the original variable. Set \(\tau(0)=0\) for the ordinary
Lebesgue integrals beginning at zero.
\end{hbdefinition}

\begin{hbtheorem}{thm:p15-feller-tail-moment-equivalence}{proved=annex:p15.feller-tail-moment-equivalence}{Exact moment identities and equivalent vanishing conditions}
For every \(t>0\),
\[
 \begin{aligned}
 \sigma_h(t)&=-\tau(t)+\frac2t\int_0^t\tau(y)\,dy,\\
 \sigma(t)&=\sigma_h(t)+t\,\mathbb P(X=-t),\\
 m(t)&=m_h(t)-t\,\mathbb P(X=-t).
 \end{aligned}
\]
Consequently the following four conditions are equivalent:
\[
 \tau(t)\to0,\qquad
 t\mathbb P(|X|>t)\to0,\qquad
 \sigma(t)\to0,\qquad
 \sigma_h(t)\to0
                         \quad(t\to\infty).
\]
The atom terms vanish in the identities when \(F\) is continuous at
\(-t\). The proof of the converse from the truncated moment uses an
explicit dyadic tail estimate; it is not left as a source referral.
\end{hbtheorem}

\section{The full source weak-law chain}
\label{sec:p15-feller-generalized-weak-law}

\begin{hbtheorem}{thm:p15-feller-truncation-inequality}{proved=annex:p15.feller-truncation-inequality}{The source truncation and Chebyshev estimate}
Let \(X_1,\ldots,X_n\) be independent with common real law \(F\),
let \(s>0\), and put
\[
 X_k'=X_k\mathbf1_{\{|X_k|\leq s\}},\quad
 S_n=\sum_{k=1}^nX_k,\quad S_n'=\sum_{k=1}^nX_k',
 \quad m_n'=n\,m(s).
\]
For every \(\varepsilon>0\),
\[
 \mathbb P(|S_n/n-m(s)|>\varepsilon)
 \leq\frac{\mathbb E[(X_1')^2]}{n\varepsilon^2}
                         +n\mathbb P(|X_1|>s).
\]
The center is the finite truncated mean, not an assumed expectation
of \(X_1\). The variance contribution has the \(n^{-1}\) scale.
\end{hbtheorem}

\begin{hbtheorem}{thm:p15-feller-generalized-weak-law}{proved=annex:p15.feller-generalized-weak-law}{Feller VII.7, Theorem 1: necessity and sufficiency}
Let \(X_1,X_2,\ldots\) be i.i.d. real random variables with law
\(F\), and write \(S_n=\sum_{k=1}^nX_k\). Real deterministic
centers \(\mu_n\) satisfying
\[
 \mathbb P(|S_n/n-\mu_n|>\varepsilon)\longrightarrow0
                           \quad\text{for every }\varepsilon>0
\]
exist if and only if \(\tau(t)\to0\) as \(t\to\infty\).
When the condition holds, one may choose
\(\mu_n=m(n)=\mathbb E[X_1\mathbf1_{\{|X_1|\leq n\}}]\);
the half-open centers \(m_h(n)\) work as well.

For necessity, let \(a\) be any median of the common law, and let
\({}^{\circ}S_n\) be the sum of its independent copied differences,
with product probability denoted by \(\mathbb Q\). The complete source chain is
\[
 \begin{aligned}
 2\mathbb P(|S_n-n\mu_n|>n\varepsilon)
 &\geq\mathbb Q(|{}^{\circ}S_n|>2n\varepsilon)\\
 &\geq\tfrac12\left(1-e^{-n\mathbb Q(|{}^{\circ}X_1|>2n\varepsilon)}\right)\\
 &\geq\tfrac12\left(1-e^{-(n/2)\mathbb P(|X_1|>2n\varepsilon+|a|)}\right).
 \end{aligned}
\]
Probabilities of copied quantities in this display are taken in their
constructed product space. The annex supplies that construction and
every step connecting the chain to \(\tau(t)\to0\).
\end{hbtheorem}

\begin{hbcorollary}{thm:p15-feller-centering-corollaries}{proved=annex:p15.feller-centering-corollaries}{All centers, finite means and symmetric laws without a mean}
Under \(\tau(t)\to0\), a deterministic sequence \(c_n\) centers
\(S_n/n\) to zero in probability exactly when \(c_n-m(n)\to0\).
If \(\mathbb E|X_1|<\infty\), then
\(\tau(t)\to0\), \(m(n)\to\mathbb EX_1\), and
\(S_n/n\to\mathbb EX_1\) in probability.
If \(F\) is symmetric, every \(m(n)=0\); hence \(\tau(t)\to0\)
gives \(S_n/n\to0\) in probability even when the first absolute
moment is infinite.
\end{hbcorollary}

\section{Worked cases and hypothesis boundaries}
\label{sec:p15-feller-worked-cases}

\begin{hbtheorem}{thm:p15-feller-symmetrization-examples}{proved=annex:p15.feller-symmetrization-examples}{The source continuous examples and an atomic calculation}
\begin{enumerate}
\item Symmetrizing an exponential law of rate \(\lambda>0\) gives
the density \((\lambda/2)e^{-\lambda|z|}\).
Symmetrizing the uniform law on \((0,1)\) gives
\((1-|z|)\mathbf1_{\{|z|\leq1\}}\).
\item If \(X\) has probabilities \(1-p,p\) at \(0,1\), where
\(0\leq p\leq1\), then \(X-X'\) has masses
\(p(1-p),p^2+(1-p)^2,p(1-p)\) at \(-1,0,1\), respectively.
In particular its CDF and its left limit at zero differ.
\item The fair law on \(\{-1,1\}\) is symmetric but cannot be
the law of \(X-X'\) for independent identically distributed real
\(X,X'\). Every such copied difference has positive probability
in each neighborhood of zero.
\item For the constant law \(X=-1\),
\[
 \tau(t)=t\mathbf1_{\{0<t\leq1\}},\quad
 \sigma(t)=t^{-1}\mathbf1_{\{t\geq1\}},\quad
 \sigma_h(t)=t^{-1}\mathbf1_{\{t>1\}}.
\]
Thus closed and half-open truncations already differ at one atom.
\end{enumerate}
\end{hbtheorem}

\begin{hbtheorem}{thm:p15-feller-proof-boundaries}{proved=annex:p15.feller-proof-boundaries}{Sharp constants, atoms and missing assumptions}
\begin{enumerate}
\item The constant two in V.(5.6) cannot be decreased uniformly:
for Bernoulli probability \(0<p<1\) and \(0<t<1\), the ratio of the
copied-difference tail to \(\mathbb P(|X|>t/2)\) is \(2(1-p)\).
The factor one half in the zero-median lower bound and in the
largest-term bound is attained by independent fair signs.
\item For independent fair signs \(Y_1\) and the constant \(Y_2=0\),
the remainder \(R\) in the largest-term construction is zero.
Its nonnegative and negative sign events do not have equal probability.
The proof's at-least-half bound still applies.
\item Dependence can invalidate the lower symmetrization bound:
take the same fair Bernoulli variable as both \(X\) and \(X'\).
Dependence can also invalidate the largest-term bound: take
\(Y_2=-Y_1\) for a fair sign. Removing symmetry invalidates that
bound even for independent constants \(Y_1=1,Y_2=-1\).
\item If \(X=1\), zero is not a median and the zero-median lower
bound would be false. For a symmetric fair sign at \(t=1\),
the source CDF form V.(5.10) cannot have its left-hand
non-strict threshold changed to a strict one.
\item For \(n=2\) independent fair signs, cutoff \(s=1\), and
\(\varepsilon=9/10\),
\(\mathbb P(|S_n/n|>\varepsilon)=1/2\).
An \(n^{-2}\) coefficient multiplying the single-summand truncated
second moment would give the false bound \(25/81\).
\item The symmetric law assigning mass
\(1/(2\log2\,j2^j)\) to each of \(\pm2^j\), \(j\geq1\),
has infinite positive and negative part expectations but satisfies
\(\tau(t)\to0\). Its i.i.d. averages tend to zero in probability.
\end{enumerate}
\end{hbtheorem}

\begin{proofboundary}{def:p15-feller-next-boundary}{What this source component does not replace}
The generalized weak law VII.7, Theorem 1 and its complete
median/symmetrization chain are supplied here. Feller's different
proof referred to in XVII.2a remains a distinct presentation.
The positive-gain normalization theorem beginning later on p.~236
and continuing onto p.~237 is supplied, with its exact normalizer
construction, in \cref{thm:p15-feller-positive-weak-normalization,thm:p15-positive-normalizer-construction}. The atomic
Poisson/Bessel specialization V.(5.5) is likewise not needed by this
chain. Those topics retain their own source and proof obligations.
\end{proofboundary}
